%% kapa: exam=1 round=33 year=2022 period=2 form=A subject=accounting subject_ko=회계학 seq=26 number=66
%% source: Q-Net 감정평가사 기출문제 (cst00309) artlSeq=5209997 file=raw/1008/5209997/2022년도 제33회 감정평가사 1차시험 2교시(감정평가 관계법규, 회계학).pdf
%% conversion: kordoc 4.24.51 -> tools/gen1.py (auto); review=proofread-sonnet; flags=none
\begin{question}{66}
무형자산의 회계처리에 관한 설명으로 옳지 않은 것은?
\begin{choices}
  \item 무형자산의 잔존가치는 해당 자산의 장부금액과 같거나 큰 금액으로 증가할 수도 있다.
  \item 브랜드, 제호, 출판표제, 고객목록, 그리고 이와 실질이 유사한 항목(외부에서 취득하였는지 또는 내부적으로 창출하였는지에 관계없이)에 대한 취득이나 완성 후의 지출은 발생시점에 항상 당기손익으로 인식한다.
  \item 무형자산의 상각방법은 자산의 경제적 효익이 소비될 것으로 예상되는 형태를 반영한 방법이어야 하지만, 그 형태를 신뢰성 있게 결정할 수 없는 경우에는 정액법을 사용한다.
  \item 내용연수가 비한정적인 무형자산은 상각하지 않고, 무형자산의 손상을 시사하는 징후가 있을 경우에 한하여 손상검사를 수행한다.
  \item 내부적으로 창출한 브랜드, 제호, 출판표제, 고객목록과 이와 실질이 유사한 항목은 무형자산으로 인식하지 아니한다.
\end{choices}
\end{question}
\answer{4}
